\scanpage{143}
% archive sheet number: 142

de $X\times_S X'$ sur $X\wedge_S X'$, et celui-ci est
le \deleted{sous} schéma quotient (universel)

\[
X\wedge_S X'\simeq X\times_S X'\big/\underbrace{\sigma_X\times_S\sigma_{X'}}_{=\sigma_P}
\]

conformément au yoga général des
produits contractés. Pour que
$X\times_S X'\unsure{-|\longrightarrow}X\wedge_S X'$ soit étale, il f. et s. qu'en
tt $s\in S$, $X_s$ \uline{ou} $X'_s$ soit étale. Pour
que $X\wedge_S X'$ soit étale sur $S$, il f et s.
que \deleted{en tout point $s\in\unsure{S|k}$}, $X$ \inserted{\uline{et}} $X'$ soient étales.

(Dans ce cas, on retrouve la notion familière.
Si $X$ \unsure{seult|seulemt|seule} est étale, on peut considérer
$X\wedge X'$ comme ``twisté'' de $X'$ par le
$(\mathbb Z/2\mathbb Z)_S$-torseur $X$ et l'op. de $\mathbb Z/2\mathbb Z$ sur $X'$
par son inv. can.)

En termes des Alg. $H,H'$ de $X,\unsure{x'|X'}$,
et $P=H\otimes H'$ de $X\times_S X'$ (avec cette fois
la structure d'Alg. comm. can. sur $P$ !]
on peut décrire celle de $X\wedge X'$ comme
la sous-Alg. \inserted{$H\wedge H'$} des invariants sous $'$.
Si $u,v$ sont des bases, alors $H\wedge H'$
a comme base $1\ (=1\otimes 1)$ et

\[
\boxed{u*v\overset{\mathrm{déf}}{=}\mathrm{Tr}(v)u+\mathrm{Tr}(u)v-2\underbrace{uv}_{\substack{\parallel\\ u\otimes v}}}
\ \in(H\otimes H')^{\unsure{\sigma_P}}
\]

% slanted
\margin{\uline{NB} Cet élément est défini sans \unsure{aucunes|aucune} conditions (\deleted{\unsure{de}} car etc) sur $H,H'$ \deleted{\unsure{ou}} \unsure{ou} $u,v$}

On a alors

